\chapter{First Contact: Functions, Variables, and the Game Loop in C}
\label{ch:c1}

This chapter begins the computer programming track. It is independent of the
assembly track: you need not have read Part~II, only the shared foundation in
Part~I. Here you will write a complete retro game in C, running on both QEMU and
the ESP32\nobreakdash-C6. The same three-function shape from Chapter~\ref{ch:platform}
organises the work, but where assembly exposed registers and the stack, C lets you
write the game's logic almost as plainly as you would describe it in words
\citep{prg32tutorialc}.

\section{Why C, and why the same platform}

C occupies a deliberate middle ground. It is a high-level language, you write
functions, name variables, and use \code{if} and \code{while}, yet it stays close
enough to the machine that nothing important is hidden. A C variable is a named
box in memory; a C function compiles to the very calling convention you would
write by hand in assembly. By learning C on the same platform, the same display,
and the same two targets used for the architecture course, you keep your attention
on the language and the game logic rather than on a new toolchain
\citep{prg32teaching}. If you later read Part~IV, you will see your C compiled into exactly the same assembly idioms used in Part~II.

\begin{prgconcept}
The platform runs C and assembly games through the identical \code{init},
\code{update}, \code{draw} contract and the identical framework interface. The
only thing that changes between the two tracks is the language you express the
logic in. This is why the platform can serve both a computer architecture course
and a computer programming course from one codebase \citep{prg32teaching}.
\end{prgconcept}

\section{The three functions, in C}

In C, the three entry points are ordinary functions returning \code{void}, named
with your game's prefix. The framework calls \code{<prefix>\_init} once and then
\code{<prefix>\_update} and \code{<prefix>\_draw} once per frame, exactly as in the
assembly track \citep{prg32tutorialc}. To use the framework, include one header.

\lstinputlisting[
    style=c,
    caption={%
        \fname{hello\_world\_c.c}: The smallest C game. Three functions that do nothing yet: the C counterpart of the assembly skeleton in Chapter~\ref{ch:asm1}.%
    }]{examples/c_first_contact/hello_world_c.c}

The single \code{\#include "prg32.h"} gives you every framework function and every
named constant: the button bits, the colours, and the screen dimensions that you
met in Chapter~\ref{ch:platform}. You never declare these yourself; the header does
it for you. For a complete list, you can check Appendix~\ref{app:api}.

\begin{prgcheck}
Build this skeleton for QEMU and confirm it runs without crashing and shows a
blank screen. As in the assembly track, an empty screen is the correct first
result: the framework has found and called your three functions.
\end{prgcheck}

\section{State: variables that persist between frames}

A game remembers things between frames. In assembly that meant the \code{.data}
section; in C it means variables declared outside the functions, at file scope,
so they keep their values from one frame to the next. The platform's C Pong
declares its paddle and ball state exactly this way \citep{prg32repo}.

\lstinputlisting[
    style=c,
    caption={%
        \fname{pong\_c\_state.c} \protect\\
       Game state as file-scope variables. Declared outside the functions, they persist across frames. \texttt{static} keeps them private to this file.}
]{examples/c_first_contact/pong_c_state.c}

The \code{init} function is where you can set the starting values of the variables.
In the pong example, this is where the initial position and velocities are set.

\lstinputlisting[
    style=c,
    caption={%
        \fname{pong\_c\_init.c} \protect\\
       Initialising state in \code{init}. This runs once, before the first frame.}
]{examples/c_first_contact/pong_c_init.c}

\begin{prgconcept}
File-scope variables in C play the role that labelled words in \code{.data} play
in assembly: they are the game's persistent state. \code{init} gives them their
opening values; \code{update} changes them each frame; \code{draw} reads them to
render. Keeping state in named variables, rather than recomputing it, is what
makes a game a game rather than a static picture.
\end{prgconcept}

\section{Reading input: the bitmask in C}

Input works just as it did in Chapter~\ref{ch:platform}: one call returns a
bitmask, and you test a button with bitwise AND against a named constant. In C the
call returns a value you store in a variable, and the test is an \code{if}.

\lstinputlisting[
    style=c,
    caption={%
        \fname{pong\_c\_update.c} \protect\\
       Reading input and moving the paddle. Note the bitwise \texttt{\&}, not \texttt{==}, exactly as warned in Chapter~\ref{ch:platform}.}
]{examples/c_first_contact/pong_c_update.c}

\begin{prgwarn}
Use \code{input \& PRG32\_BTN\_LEFT}, with a single ampersand for bitwise AND, not
\code{input == PRG32\_BTN\_LEFT}. The mask may hold several buttons at once; the
bitwise test asks ``is the left bit set, regardless of the others,'' which is what
you want. This is the same warning as in the assembly track, and it catches
beginners in both languages \citep{prg32tutorialc}.
\end{prgwarn}

The two clamps at the end are the boundary check from Chapter~\ref{ch:asm2},
written in C. The maximum is the screen width minus the paddle's own width, so the
paddle's right edge never leaves the viewport. In C the clamp is two short
\code{if} statements; the logic is identical to the assembly version, only more
compact.

\section{The ball: velocity and bounce, in C}

The ball moves by adding its velocity to its position each frame and bounces by
negating its velocity at the edges: the same rule as in
Chapter~\ref{ch:asm3}, now in three lines of C \citep{prg32repo}.

\lstinputlisting[
    style=c,
    caption={%
        \fname{pong\_c\_ball\_motion.c} \protect\\
        Ball motion and edge bounce. Adding velocity moves the ball; negating it at the boundary turns it around.}
]{examples/c_first_contact/pong_c_ball_motion.c}

\section{Collision: one call, one decision}

To bounce the ball off the paddle, ask the framework whether their rectangles
overlap. \code{prg32\_sprite\_hitbox} takes the position and size of two
rectangles and returns non-zero when they touch, so collision is a single call
inside an \code{if} \citep{prg32repo}.

\lstinputlisting[
    style=c,
    caption={%
        \fname{pong\_c\_paddle\_collision.c} \protect\\
        Bouncing the ball off the paddle. The hitbox test returns non-zero on overlap.}
]{examples/c_first_contact/pong_c_paddle_collision.c}

\begin{prgconcept}
Collision detection in C is one function call and one \code{if}. The framework
does the geometric work of testing overlap on both axes; you decide what happens
when the answer is yes. This separation, the engine answers ``did they touch,''
your code decides ``so what'', is a pattern you will meet throughout game
programming.
\end{prgconcept}

\section{Drawing the frame}

The \code{draw} function clears the viewport and renders each object as a coloured
rectangle, using the named colour constants from the header. The \code{draw} function 
reads the state variables but does not change them; all logic stays in \code{update}.
We do this to follow the phase separation from Chapter~\ref{ch:platform}.
\citep{prg32tutorialc}.

\lstinputlisting[
    style=c,
    caption={%
        \fname{pong\_c\_draw.c} \protect\\
        Rendering the frame. \code{draw} only reads state and puts pixels; it changes nothing.}
]{examples/c_first_contact/pong_c_draw.c}

Now we can look at the full Pong example, built in C. If you followed the assembly course,
you can compare it line by line with the previous version.

\lstinputlisting[
    style=c,
    caption={%
        \fname{pong\_skeleton.c} \protect\\
        Rendering the frame. \code{draw} only reads state and puts pixels; it changes nothing.}
]{examples/c_first_contact/pong_skeleton.c}

\begin{prgtargets}
\textbf{Both targets.} This complete C Pong runs unchanged on QEMU and the
ESP32\nobreakdash-C6. On the emulator you watch the ball bounce and verify the layout; on
the board you play it with a real joystick.
\end{prgtargets}

\begin{prgtry}
Change the paddle colour from \code{PRG32\_COLOR\_WHITE} to
\code{PRG32\_COLOR\_CYAN} and the ball size from \code{8} to \code{12}, then
rebuild. Both changes should be visible immediately. Next, the platform's
break-and-fix exercise: delete the \code{paddle\_x > 256} clamp, rebuild, hold
RIGHT until the paddle leaves the screen, then restore the clamp and explain in one
sentence which condition protects the boundary \citep{prg32tutorialc}.
\end{prgtry}

\begin{prgcheck}
You can now write a complete C game: persistent state in file-scope variables,
input read as a bitmask and tested with \code{\&}, motion by velocity with
boundary clamps, bounce by negation, collision by one call, and rendering kept
separate in \code{draw}. Chapter~\ref{ch:c2} scales these ideas up with the C
features, arrays and structs, that let a game hold many objects and a whole
world.
\end{prgcheck}
